% RESEARCH_DRAFT_SHA256: 3c09eafe0ab851b6554c3d7ea970d6963d39b3a492f3f9eeca506a95eb77d730
% Layout and source-reference wrappers expanded for the web renderer.
\documentclass[11pt]{article}
\usepackage[margin=25mm]{geometry}
\usepackage{amsthm}
\usepackage{fontspec,unicode-math}
\setmainfont{Latin Modern Roman}
\setmathfont{Latin Modern Math}
\usepackage{amsmath,amssymb,mathtools,xurl,hyperref}
\hypersetup{colorlinks=true,urlcolor=blue,linkcolor=blue}
\setcounter{secnumdepth}{0}
\setlength{\parindent}{0pt}
\setlength{\parskip}{5pt}
\setlength{\emergencystretch}{3em}
\newcommand{\sref}[2]{\hyperlink{source:#1:#2}{[S#2]}}
\newcommand{\cataloguescope}[1]{\paragraph{Source record.}#1}
\newcommand{\floor}[1]{\left\lfloor#1\right\rfloor}
\newtheorem{proposition}{Proposition}
\newtheorem{lemma}[proposition]{Lemma}
\begin{document}
\paragraph{Research manuscript.} OpenAI \texttt{gpt-6-astra}, reasoning \texttt{ultra}; 11 October 2026.
\paragraph{Preserved source body.} The complete catalogue mathematical text and sources follow with layout and source links expanded, including source qualifications and its historical preparation notes. New research begins at the section headed Research attempt.

\section{Restricted Models of Vertex-Distribution-Free Testing}\label{statement}

\subsection{Definitions and mathematical statement}

\paragraph{Shared notation.}Standard number systems are $\mathbb N,\mathbb Z,\mathbb Q,\mathbb R,\mathbb C$. Specialized notation and index ranges follow the cited source and the definitions below.


In vertex-distribution-free testing, vertices are sampled from an unknown distribution. Distance weights changed adjacency pairs by the product of their vertex probabilities. Restricting atom weights limits how heavily or lightly individual vertices may be sampled.
\paragraph{Record scope and qualifications.}
The catalogue title is retained as metadata. Named certificates, contracts, bounds or frameworks in that title are not asserted as results of the original source extract.
The primary report states theorems for the named restricted directions. The cleaned question is retained as framing and does not identify a separate remaining theorem beyond those results.

Source-record limitation: The primary report immediately states theorems resolving the three restricted-model directions (large graphs, small maximum atoms, and lower atom bounds). The supplied question is a broad framing paragraph and does not identify a remaining mathematical target beyond those stated results.
\paragraph{Statement recorded in the supplied source.}
How does vertex-distribution-free graph-property testing change if testers are required to work only on sufficiently large graphs, or if the vertex distribution is restricted by lower or upper bounds on its atom weights? \href{https://doi.org/10.4171/owr/2019/19}{[S2]}
\paragraph{Source status.}
The supplied catalogue labels this source record open; that metadata was checked against the supplied snapshot on 2026-10-11. The label is preserved for provenance and is not an independent verification of present open status or mathematical correctness. The catalogue literature assessment is dated 2026-08-22.
\subsection{Short English statement}
Study graph testing under restricted vertex distributions and graph sizes.
\subsection{Sources}
\begin{itemize}
\item[S1] Ulam AI and the UnsolvedMath Contributors, UnsolvedMath: A Curated Collection of Open Mathematics Problems, catalogue entry. Catalogue curation and metadata: CC BY 4.0. The source extract is retained; definitions, English summaries and qualifications are editorial additions.
\url{https://huggingface.co/datasets/ulamai/UnsolvedMath}.
\item[S2] Combinatorics, Probability and Computing (2020), Oberwolfach Reports, DOI 10.4171/owr/2019/19, pp. 1113--1189
\url{https://doi.org/10.4171/owr/2019/19}.
\item[S3] Goldreich--Orenstein, Testing properties in the vertex-distribution-free model
\url{https://eccc.weizmann.ac.il/report/2019/083/revision/1/download/}.
\item Source-specific Oberwolfach Reports licences, including CC BY-SA 4.0 where applicable, are retained; see THIRD\_PARTY\_NOTICES.md.
\end{itemize}
\paragraph{Preparation and review.}
The original source and UnsolvedMath contributors retain attribution for the source extract. Codex (GPT-6; exact serving revision not exposed) assisted with definitions, English summaries, scope notes and formatting. No new human mathematical authorship or independently verified solution is asserted. Catalogue identity checked against the complete supplied registry and current repository. Distinct upstream records, including repeated questions and broad or damaged extracts, are retained for complete coverage of the source CS category. A repeated record is not a new independent conjecture. Complete mathematical-equivalence and current-literature review has not been performed. Preparation and final review follow contribution prompts C and D, with source gaps disclosed. Structural and rendering checks do not establish a new theorem or open status.
\label{end:statement}

\section{Research attempt}
\subsection{Literature, source context and scope}
The full original catalogue body and source qualifications are preserved above. The relevant primary workshop passages were read in the cached publisher report. Historical claims and subsequent attributed results are distinguished from the mathematical work below. A restricted example, standard lemma or known theorem is not counted as new progress on the entire original target. Missing assumptions and repeated or rhetorical extracts are not silently turned into different solved problems.

\subsection{Atom bounds compare distances but do not give the missing oracle}
The preserved source limitation is substantive: the report already states its large-graph, small-maximum-atom and lower-atom theorems. We do not relabel that framing paragraph as a new unresolved classification. Instead we test a simple quantitative transfer under a stronger, explicit hypothesis. On an $n$-vertex set suppose
\[
\frac c n\le D(v)\le\frac C n\quad\hbox{for all }v,
\qquad 0<c\le1\le C.
\]
For a graph property admitting at least one graph on this vertex set, let $d_D(G,\mathcal P)$ denote the minimum weighted unordered-edge edit distance, and let $d_U$ use the uniform distribution.
\begin{proposition}
One has $c^2d_U(G,\mathcal P)\le d_D(G,\mathcal P)\le C^2d_U(G,\mathcal P)$. If an additional oracle gives exact values of $D(v)$ and $n$ is known, rejection sampling turns $D$-samples into uniform samples with expected $1/c$ trials per accepted sample.
\end{proposition}
\begin{proof}
Every changed unordered pair has weight between $c^2/n^2$ and $C^2/n^2$. Summing proves the two inequalities for each candidate edited graph. Taking minima proves the lower inequality, and choosing a uniform-distance minimizer proves the upper inequality; the finite set of graphs guarantees a minimizer. For rejection sampling, accept a sample $v$ with probability $c/(nD(v))$, which lies in $[0,1]$. The joint probability of drawing and accepting $v$ is exactly $c/n$, hence the total success probability is $c$ and its conditional distribution is uniform. Independent trials give mean waiting time $1/c$ by summing the geometric tail.
\end{proof}
This yields an explicit attempted transfer: $\varepsilon$-far in $D$ implies $\varepsilon/C^2$-far in uniform distance, then simulate a uniform tester by rejection. The first step is valid; the second requires extra information. The actual VDF model supplies samples from an unknown $D$, not exact atom evaluations, and it does not provide $n$. The displayed simulation is therefore not an algorithm in that model. Moreover an expected sample count alone is not the source's always-bounded query guarantee without a separate truncation/error analysis.
\subsection{Verification and remaining gaps}
For $D=(1/2,1/4,1/4)$ on three vertices, the bounds with $c=3/4$ and $C=3/2$ were checked on every changed-edge subset, and the accepted distribution was computed exactly. These algebraic comparisons help separate distance control from access assumptions. They provide no new proof of the source's already-stated restricted-model theorems, no query-optimal tester, and no definition of an additional open target hidden in its framing paragraph.

\section{References and source attribution}
Primary workshop report: \url{https://doi.org/10.4171/owr/2019/19}. Relevant passages read: 2019/19 pp. 1138--1140, unknown-distribution access model and Theorems6--8.

All retrievals were checked on 11 October 2026. Original question attribution remains in the preserved Sources section. External theorems are attributed and not claimed as new discoveries or fully verified merely by citation.


\paragraph{Integration review.} Codex (GPT-6; exact serving revision not exposed) reviewed the written arguments and expanded layout and source-reference wrappers. Repository metadata and the production web renderer were checked using Python 3.13 and JavaScript with MathJax 3.2.2. This model-assisted review is not independent human proof verification; no PDF compilation or formal proof checking is asserted.

\section{Completion Estimate}
\noindent\textbf{COMPLETION: 0\%}
No substantiated new progress toward the entire original target is claimed. The worked results and failed route are delimited above; a known result, illustrative model, or auxiliary obstruction is not counted as original-target completion.

\section{Final status and attribution}
\textbf{Attempt status: unresolved.} This describes this research attempt, not a new ruling on the source catalogue's status or on external literature claims.
The mathematics and this self-audit were generated by OpenAI \texttt{gpt-6-astra}, reasoning \texttt{ultra}, on 11 October 2026. Codex (GPT-6; exact serving revision not exposed) orchestrated the batch, reviewed the written mathematical arguments, and contributed formatting/repository review. This additional model review is not independent human verification. No human mathematical authorship, independent human proof checking, or formal verification is claimed. The historical preparation paragraphs in the reproduced catalogue body concern that body only.

Contribution prompt A was applied to the complete definition and supplied source record, followed by the adversarial checks of prompt B. Local reproduction files are \texttt{batch022/checks.py}, \texttt{computations.json}, and the per-ID audit JSON. They were not public links at preparation time. Numerical checks are finite and do not prove unrestricted results. No LaTeX engine was found locally; compilation was not performed by this research agent.

\end{document}
